\begin{titlepage}

\thispagestyle{empty}
\vspace{-2cm}

\voffset-10pt


%%%
%% premiere ligne avec logo légaux et numéro national de thèse
%% (modifier les symboles % en début de lignes 
%%  pour obtenir la configuration souhaitée)
%%%
\noindent
\hspace*{-1cm}\hbox{\includegraphics[width=8.6cm]{images/ed_edmh-h.jpg}} 
\hfill 
\hspace*{-0.4cm}{\small {\bf NNT : 2016SACLY000}}
%%%
%% Numéro National de Thèse, à préciser lors du second dépôt 
%% de la thèse post-soutenance
%%%
\hfill
%%%
%% logo Etablissement inscripteur
%% (modifier les symboles % en début de lignes 
%% pour obtenir la configuration souhaitée, et ajuster 
%% éventuellement la taille) En cas de cotutelle internationale, 
%% mettre à la place le logo de l'université étrangère et rajouter 
%% en bas de page le logo de l'établissement inscripteur
%%%
%\hbox{\includegraphics[width=2.3cm]{logoCentraleSupelec.jpg}}
%\hbox{\includegraphics[width=2cm]{logoHEC.jpg}}
%\hbox{\includegraphics[width=2cm]{logoENSAE.jpg}}
%\hbox{\includegraphics[width=2cm]{logoENSCachan.jpg}}
%\hbox{\includegraphics[width=1.5cm]{logoENS.jpg}}
%\hbox{\includegraphics[width=1.9cm]{logoENSTA.jpg}}
\hbox{\includegraphics[width=2.7cm]{logoUPSud_UPSaclay.jpg}}
%\hbox{\includegraphics[width=2cm]{logoTelecom.jpg}}
%\hbox{\includegraphics[width=2cm]{logoEvry.jpg}}
%\hbox{\includegraphics[width=2.2cm]{logoUVSQ.jpg}}
%\hbox{\includegraphics[width=2.5cm]{logoX.jpg}}
\vspace{7mm}

\begin{center}
{\Large\bf TH\`ESE DE DOCTORAT}
\end{center}
\begin{center}
{de }
\end{center}
\begin{center}
 {\Large\sc l'Universit\'e Paris-Saclay}\\
  \vspace*{0.4cm}
\'Ecole doctorale de math\'ematiques Hadamard (EDMH, ED 574)\\  
 \vspace*{0.4cm} 
%%%
%% nom \'etablissement inscripteur
%%%
{\small \it \'Etablissement d'inscription : } 
%    Centrale-Sup\'elec\\
%    Ecole polytechnique\\
%    Ecole normale sup\'erieure\\
%    Ecole normale sup\'erieure de Cachan\\
%    \'Ecole nationale de la statistique et de l'administration \'economique\\
%    \'Ecole nationale sup\'erieure de techniques avanc\'ees\\
%    Ecole des hautes \'etudes commerciales de Paris\\
%    Telecom SudParis\\
%    Telecom ParisTech\\
%    Universit\'e d'Evry-Val d'Essonne\\
    Universit\'e Paris-Sud\\
%   Universit\'e de Versailles Saint-Quentin-en-Yvelines\\  
 \vspace*{0.2cm} 
%%%
%% nom \'etablissement(s) d'accueil, si diff\'erent du pr\'ec\'edent
%%%
{\small \it \'Etablissement d'accueil : }    
%    AgroParisTech\\
%    Commissariat \`a l'\'energie atomique et aux \'energies alternatives\\
    Institut des hautes \'etudes scientifiques\\
%    Institut national de la recherche agronomique\\
\vspace*{0.2cm} 
%%%
%% nom laboratoire(s) d'accueil
%%%
{\small \it Laboratoire d'accueil : }     
%  Centre de math\'ematiques appliqu\'ees de Polytechnique, UMR 7641 CNRS\\
%  Centre de math\'ematiques et de leurs applications, UMR 8536 CNRS\\
%  Centre de math\'ematiques Laurent Schwartz, UMR 7640 CNRS\\
%  D\'epartement de math\'ematiques et applications, UMR 8553 CNRS\\
%  ENSAE-X Centre d'\'economie, statistique et sociologie, UMR 9194 CNRS\\
%  F\'ed\'eration de Math\'ematiques, FR 3487 CNRS\\
%  Groupement de recherche et d'\'etudes en gestion (GREGHEC), UMR 8071 CNRS\\
%  Institut de physique th\'eorique de Saclay, URA 2306 CNRS\\
Laboratoire Alexander Grothendieck (ERL 9216 CNRS)\\
%  Laboratoire de math\'ematiques d'Orsay, UMR 8628 CNRS\\
%  Laboratoire de math\'ematiques de Versailles, UMR 8100 CNRS\\
%  Laboratoire de math\'ematiques et mod\'elisation d’\'Evry, UMR 8071 CNRS-INRA\\
%  Laboratoire traitement et communication de l'information, UMR 5141 CNRS\\
%  Math\'ematiques et informatique appliqu\'ees, UMR 518 INRA\\
%  Math\'ematiques et informatique appliqu\'ees du g\'enome \`a l'environnement, UR 1404
%  Services r\'epartis, architectures, mod\'elisation, validation, administration des r\'eseaux, UMR 5157 CNRS\\
%  Unit\'e de math\'ematiques appliqu\'ees, ENSTA-CNRS-INRIA\\
\vspace*{0.2cm}
\end{center}




\begin{center}
{\it Sp\'ecialit\'e de doctorat : } 
%%%
%%  Choisir l'une des trois spécialités suivantes (soumis à
%%  accord du comité de direction de l'EDMH)
%%%
{\large Math\'ematiques fondamentales}
%{\large Math\'ematiques appliquées}
%{\large Math\'ematiques aux interfaces}
\end{center}

\vspace{5mm}

\begin{center}
{\large\bf Pr\'enom NOM}
\end{center}

\vspace{3mm}
 
\begin{center}
{\Large Titre de la th\`ese}
%%%
%% éventuellement sur deux lignes
%%%
\end{center}


\vspace{10mm}

\noindent{\small \it Date de soutenance~: } Jour Mois Ann\'ee

\vspace{5mm}

\noindent
{\small \it Apr\`es avis des rapporteurs~: }
\begin{tabular}{l} {\sc Pr\'enom NOM} (Institution)\vspace{1mm}  \\
{\sc  Pr\'enom NOM} (Institution)\\
\end{tabular}

\vspace{8mm}

\noindent
{\small \it Jury de soutenance~: }
%%%
%% par ordre alphabétique des membres
%%
%% version avant soutenance à adapter (un jury peut contenir 
%% tout ou partie des rapporteurs, tout au partie des codirecteurs de thèses,
%% voire des invités). Le président du jury est choisi par le jury en 
%% son sein le jour de la soutenance, et est indiqué sur les exemplaires 
%% post-soutenance
%% En cas d'absense d'un membre de jury prévu, les exemplaires post-thèse 
%% ne doivent faire figurer que les présents.
%%%%
\begin{tabular}{ll}
{\sc   Pr\'enom NOM}&(Institution) {\small Rapporteur}\vspace{1mm}\\
{\sc   Pr\'enom NOM}&(Institution) {\small Codirecteur de th\`ese}\vspace{1mm}\\
{\sc   Pr\'enom NOM}&(Institution) {\small Examinateur}\vspace{1mm}\\
{\sc   Pr\'enom NOM}&(Institution) {\small Examinateur}\vspace{1mm}\\
{\sc   Pr\'enom NOM}&(Institution) {\small Codirecteur de th\`ese}\vspace{1mm}\\
{\sc   Pr\'enom NOM}&(Institution) {\small Invit\'e}\vspace{1mm}\\
\end{tabular}


\vfill
\noindent
\hbox{\includegraphics[width=2.2cm]{logo_fmjh.jpg}}
\hfill 
%%%
%% logo Etablissement d'accueil, si différent du précédent
%% (modifier les symboles % en début de lignes 
%% pour obtenir la configuration souhaitée, et ajuster 
%% éventuellement la taille)
%%%
\hbox{\includegraphics[width=2cm]{logoIHES.jpg}}
%\hbox{\includegraphics[width=2.1cm]{logoINRA.jpg}}
%\hbox{\includegraphics[width=2.5cm]{logoAgro.jpg}}
%\hbox{\includegraphics[width=1.5cm]{logoCEA.jpg}}
\hfill
%%%
%% logo laboratoire(s) d'accueil, s'il existe
%%%
%\hbox{\includegraphics[width=2cm]{logoCMAP.jpg}}
%\hbox{\includegraphics[width=2.4cm]{logoCMLA.jpg}}
%\hbox{\includegraphics[width=2cm]{logoCMLS.jpg}}
%\hbox{\includegraphics[width=4cm]{logoDMA.jpeg}}
%\hbox{\includegraphics[width=2.2cm]{logoGREHEC.jpg}}
%\hbox{\includegraphics[width=2.1cm]{logoIPhT.jpeg}}
\hbox{\includegraphics[width=2cm]{logoLAG.jpg}}
%\hbox{\includegraphics[width=1.6cm]{logoLAMME.jpg}}
%\hbox{\includegraphics[width=3cm]{logoLMO.jpg}}
%\hbox{\includegraphics[width=3.4cm]{logoLMV.jpeg}}
%\hbox{\includegraphics[width=2.7cm]{logoMaIAGE.png}}
%\hbox{\includegraphics[width=2.5cm]{logoSAMOVAR.jpg}}
%\hbox{\includegraphics[width=2cm]{logoUMA.jpg}}
\hfill \includegraphics[width=1cm]{pictoParis-Saclay.jpg}




\newpage
%%%
%%
%% Quatrième de couverture (last cover page of the memoirs)
%%
%%%

\hbox{\includegraphics[width=8.6cm]{ED_EDMH-h.jpg}}


\bigskip
\noindent\fbox{\parbox{\textwidth}{
{\bf Titre : }  Titre en français

\medskip
{\bf Mots Clefs : }  Mettre de 3 à 6 mots clefs

\medskip
{\bf Résumé : } 


\vspace{7cm}

}}

\bigskip
\noindent\fbox{\parbox{\textwidth}{
{\bf Title : }  Title in english

\medskip
{\bf Keys words : }  3 to 6  key words (in english)

\medskip
{\bf Abstract : } 


\vspace{7cm}

}}

\vfill
\hfill \includegraphics[width=1cm]{pictoParis-Saclay.jpg}


\end{titlepage}
\sloppy
\titlepage
